\subsection{Case of coverings associated with a principal covering}

Let B be a topological space, let G be a discrete topological group, and let E be a principal covering of B with group G. Denote by p the projection of the B-space E. Let b be a point of B and let x be a point of \(E_{b}\) .

PROPOSITION 5. --- The map \(h_{(\mathrm{E},x)}\) from \(\pi_1(\mathrm{B},b)\) to G which, to \(\gamma \in \pi_1(\mathrm{B},b)\) , associates the unique element \(g\) of G such that \(x \cdot g = x \cdot \gamma^{-1}\) is a group homomorphism, whose kernel is the subgroup \(p_*(\pi_1(\mathrm{E},x))\) of \(\pi_1(\mathrm{B},b)\) . If E is connected, this homomorphism is surjective.

For every \(g \in \mathrm{G}\) , we have \(h_{(\mathrm{E},x \cdot g)} = \operatorname{Int}(g^{-1}) \circ h_{(\mathrm{E},x)}\) .

Let E be a principal covering of B with group G, let x be a point of \(E_{b}\) and denote by p the projection of E. For every \(g \in G\) , the map \(y \mapsto y \cdot g\) is a B-automorphism of E. We therefore have, for every \(g \in G\) , every \(y \in E_{b}\) and every \(\gamma \in \pi_{1}(B, b)\) , the relation \((y \cdot g) \cdot \gamma = (y \cdot \gamma) \cdot g\) (cf. III, p. 305, relation (4)). Consequently, for \(\gamma, \delta \in \pi_{1}(B, b)\) , we have

\[
\begin{array}{r l} & {x \cdot h _ {(\mathrm{E}, x)} (\gamma \delta) = x \cdot \delta^ {- 1} \gamma^ {- 1}} \\ & {\qquad = (x \cdot h _ {(\mathrm{E}, x)} (\delta)) \cdot \gamma^ {- 1}} \\ & {\qquad = (x \cdot \gamma^ {- 1}) \cdot h _ {(\mathrm{E}, x)} (\delta)} \\ & {\qquad = x \cdot h _ {(\mathrm{E}, x)} (\gamma) h _ {(\mathrm{E}, x)} (\delta),} \end{array}
\]

which proves that \(h_{(\mathrm{E},x)}\) is a group homomorphism.

Its kernel is the stabilizer of \(x\) for the canonical action of \(\pi_1(\mathrm{B}, b)\), and therefore is equal to \(p_*(\pi_1(\mathrm{E}, x))\) by Theorem 1 of III, p. 305. The map \(g \mapsto x \cdot g\) is a bijection from G onto \(\mathrm{E}_b\). The image of the homomorphism \(h_{(\mathrm{E}, x)}\) is therefore the set of \(g \in \mathrm{G}\) such that \(x \cdot g\) belongs to the orbit of \(x\) for this operation. If E is connected, \(\pi_1(\mathrm{B}, b)\) acts transitively on \(\mathrm{E}_b\) (loc. cit.) and the homomorphism \(h_{(\mathrm{E}, x)}\) is surjective.

Let \(g\in \mathbf{G}\) ; for every \(\gamma \in \pi_1(\mathrm{B},b)\) , we have

\[
x \cdot g h _ {(\mathrm{E}, x \cdot g)} (\gamma) = (x \cdot g) \cdot \gamma^ {- 1} = (x \cdot \gamma^ {- 1}) \cdot g = x \cdot h _ {(\mathrm{E}, x)} (\gamma) g,
\]

whence \(h_{(\mathrm{E},x\cdot g)}(\gamma)=g^{-1}h_{(\mathrm{E},x)}(\gamma)g\) . This completes the proof of the proposition.

Examples. --- 1) Let F be a discrete set equipped with an action of G and let \(X = E \times^{G} F\) be the covering of B associated with it. Denote by \(\varphi: E \times F \to X\) the canonical B-morphism. It is a morphism of coverings. For every \(\gamma \in \pi_{1}(B, b)\) and every \(f \in F\) , we then have

\[
\varphi (x, f) \cdot \gamma = \varphi (x \cdot \gamma , f) = \varphi (x \cdot h _ {(\mathrm{E}, x)} (\gamma) ^ {- 1}, f) = \varphi (x, h _ {(\mathrm{E}, x)} (\gamma^ {- 1}) \cdot f).\tag{5}
\]

If one identifies F with \(X_{b}\) by the bijective map \(f \mapsto \varphi(x, f)\) , it follows that the right operation \(\pi_{1}(B, b) \to \operatorname{Aut}(X_{b})^{\circ}\) is the composition of the homomorphism \(h_{(E,x)}\) , the antihomomorphism \(g \mapsto g^{-1}\) of G into itself and the homomorphism G → Aut(F) induced by the action of G on F.

\begin{enumerate}
\def\labelenumi{\arabic{enumi})}
\setcounter{enumi}{1}
\tightlist
\item
  Let H be a discrete topological group, let \(f\colon G\to H\) be a group homomorphism. Endow H with the left action of G given by \(g\cdot h=f(g)h\) , for \(g\in G\) and \(h\in H\) . Let \(\mathrm{E}^{\prime}=\mathrm{E}\times^{\mathrm{G}}\mathrm{H}\) be the associated covering; it is a principal covering with group H (I, p. 107, example 6). Denote by \(q\colon E\times H\to E\times^{G}H\) the canonical map and set \(x^{\prime}=q(x,e)\) . We have \(h_{(E^{\prime},x^{\prime})}=f\circ h_{(E,x)}\) .
\end{enumerate}

Indeed, let \(c\) be a loop at \(b\) in B, let \(\gamma\) be its class in \(\pi_1(B,b)\) and let \(g = h_{(E,x)}(\gamma)\) ; one therefore has \(x \cdot \gamma = x \cdot g^{-1}\) . Let \(\widetilde{c}\) be a path with origin \(x\) in E; then, \(t \mapsto q(\widetilde{c}(t), e)\) is a path with origin \(x'\) in E' which lifts the path \(p \circ \widetilde{c}\) in B, so that

\[
x ^ {\prime} \cdot \gamma = q (x \cdot \gamma , e) = q (x \cdot g ^ {- 1}, e) = q (x, f (g) ^ {- 1}) = q (x, e) \cdot f (g) ^ {- 1},
\]

which proves that \(h_{(\mathrm{E}', x')}(\gamma) = f(g)\) .

